\section*{Problem 5: Reed--Solomon Encoding (Level 5)}

\subsubsection*{Mathematical part}
\textbf{Recall:} We now possess everything required to construct a practical error-correcting code. How can we encode a message so that transmission errors become detectable and correctable?

Instead of storing the message polynomial directly, we force the encoded codeword $c(x)$ to vanish at a predetermined set of points:
\[
c(\alpha^0)=0, \quad c(\alpha^1)=0, \quad \ldots, \quad c(\alpha^{2t-1})=0
\]
where $t$ is the number of errors we wish to correct, requiring $2t$ parity symbols.

To guarantee that $c(x)$ has these roots, we construct a \textbf{generator polynomial} $g(x)$ whose roots are exactly these points:
\[
g(x) = (x-\alpha^0)(x-\alpha^1)\cdots(x-\alpha^{2t-1})
\]
Any polynomial divisible by $g(x)$ will automatically vanish at all these roots.

\paragraph{How do we verify $g(x)$?} We can check our generator polynomial by simply evaluating it at each expected root. If the construction is correct, every evaluation must yield zero:
\[
g(\alpha^0) \stackrel{?}{=} 0, \quad g(\alpha^1) \stackrel{?}{=} 0, \quad \ldots, \quad g(\alpha^{2t-1}) \stackrel{?}{=} 0
\]

\paragraph{Systematic Encoding}

The transmitted codeword is $(m_{k-1}, \ldots, m_0, p_{2t-1}, \ldots, p_0)$. \textit{Notice that the ISO QR Code standard transmits coefficients from High-to-Low degree. This ensures the text message characters are read first in natural left-to-right order, followed by the parity bytes. Our interactive UI dashboard strictly follows this High-to-Low ISO standard for all visual displays, even though your underlying Python engine computes everything Low-to-High.}

To encode a message $m(x)$ of degree $k-1$ into a codeword $c(x)$ that is divisible by $g(x)$, we use \textbf{systematic encoding}:
\begin{enumerate}
    \item \textbf{Shift:} Multiply $m(x)$ by $x^{2t}$ to make room for the parity bytes at the low-degree positions.
    \item \textbf{Divide:} Compute the remainder $r(x) = m(x) \cdot x^{2t} \pmod{g(x)}$ using polynomial long division.
    \item \textbf{Subtract:} The codeword is $c(x) = m(x) \cdot x^{2t} - r(x)$.
\end{enumerate}
Since $m(x) \cdot x^{2t} = q(x) \cdot g(x) + r(x)$, subtracting the remainder gives $c(x) = q(x) \cdot g(x)$. The codeword is divisible by $g(x)$ by construction! The encoding reuses your polynomial \lstinline|__divmod__| from Level~2.

\paragraph{Worked Example}
Suppose we wish to correct $t=2$ errors, requiring $2t = 4$ parity symbols. The generator polynomial is:
\[
g(x) = (x-\alpha^0)(x-\alpha^1)(x-\alpha^2)(x-\alpha^3)
\]
If our original message is $m(x)$, systematic encoding first multiplies it by $x^{4}$ to make room for the parity symbols. Then, we find the remainder $r(x) = m(x)x^4 \pmod{g(x)}$. The final codeword transmitted is:
\[
c(x) = m(x)x^4 - r(x)
\]
By polynomial arithmetic, $c(x)$ is perfectly divisible by $g(x)$, guaranteeing that evaluating $c(x)$ at $\alpha^0, \dots, \alpha^3$ will yield zero.

\subsubsection*{Implementation part}
\textbf{Tasks:} Implement \lstinline|get_generator_poly| and \lstinline|rs_encode|.

\textbf{Details for get\_generator\_poly:} 
You must construct the generator polynomial by iterative multiplication.
\begin{itemize}
    \item Start with the identity polynomial $g(x) = 1$.
    \item Loop through $i$ from $0$ to $2t-1$ (where $2t$ is \lstinline|num_ec_bytes|).
    \item Multiply $g(x)$ by the linear factor $(x - \alpha^i)$. In our Low-to-High coefficient ordering, this translates to \lstinline|Polynomial([-(field.alpha ** i), field.one])|. Your \lstinline|__neg__| and \lstinline|__pow__| from Level~1 handle the negation and exponentiation seamlessly.
\end{itemize}

\textbf{Details for rs\_encode:}
You must perform systematic encoding of the message polynomial.
\begin{itemize}
    \item First, get the generator polynomial $g(x)$ using \lstinline|get_generator_poly|.
    \item Construct a shift polynomial $x^{2t}$. In Low-to-High order, this is a list of \lstinline|num_ec_bytes| zeros followed by a one: \lstinline|Polynomial([field.zero] * num_ec_bytes + [field.one])|.
    \item Multiply the message polynomial by the shift polynomial to make room for the parity bytes.
    \item Compute the remainder using polynomial \lstinline|divmod| with the generator polynomial.
    \item Return the shifted polynomial minus the remainder.
\end{itemize}

\begin{center}
\renewcommand{\arraystretch}{1.5}
\begin{tabularx}{\textwidth}{@{} >{\bfseries}l X c @{}}
\toprule
Function & Arguments & Returns \\ \midrule
get\_generator\_poly & \lstinline|num_ec_bytes: int, field: ExtensionField| & \lstinline|Polynomial| \\
rs\_encode & \lstinline|msg_poly: Polynomial, num_ec_bytes: int, field: ExtensionField| & \lstinline|Polynomial| \\
\bottomrule
\end{tabularx}
\end{center}

\subsubsection*{Experimentation part}
\textbf{UI Guide:} In the interactive dashboard, explore the RS Encoding tab. The UI displays a \textbf{Root Verification Table} that evaluates your generator polynomial at each expected root $\alpha^0, \ldots, \alpha^{2t-1}$. Every row must show zero --- this is a self-contained proof that your polynomial is correct, without needing to encode anything. Below the table, you can input text to watch the systematic encoding produce a codeword with message and parity bytes.

\vspace{1cm}
